\documentclass[handout]{beamer}

\usepackage[utf8]{inputenc}
\usepackage[T1]{fontenc}
\usepackage[brazil]{babel}
\usepackage{listings}
\usepackage{mathtools}
\usepackage{amssymb}

\usetheme{Madrid}

\title{\LaTeX\ para gente grande}
\subtitle{Dia 2: Documentos Técnicos e Escrita Matemática}
\author[Andrês Oliveira]{Andrês Oliveira}
\date{Minicurso LEM/UNICAMP}

\definecolor{codegreen}{rgb}{0,0.6,0}
\definecolor{codegray}{rgb}{0.5,0.5,0.5}
\definecolor{codepurple}{rgb}{0.58,0,0.82}
\definecolor{backcolour}{rgb}{0.95,0.95,0.92}

\lstset{
    backgroundcolor=\color{backcolour},
    commentstyle=\color{codegreen},
    keywordstyle=\color{magenta},
    numberstyle=\tiny\color{codegray},
    stringstyle=\color{codepurple},
    basicstyle=\ttfamily\scriptsize,
    breaklines=true,
    showstringspaces=false,
    tabsize=2
}

\DeclareMathOperator{\rank}{rank}
\DeclareMathOperator{\tr}{tr}
\DeclareMathOperator*{\argmin}{arg\,min}
\DeclarePairedDelimiter{\abs}{\lvert}{\rvert}
\DeclarePairedDelimiter{\norm}{\lVert}{\rVert}
\DeclarePairedDelimiterX{\inner}[2]{\langle}{\rangle}{#1,#2}

\begin{document}

% BLOCO 1 - ABERTURA
\begin{frame}
    \maketitle
\end{frame}

\begin{frame}{Objetivos da Tarde}
    \begin{itemize}
        \item Organizar um documento técnico em arquivos coerentes.
        \item Escolher ambientes matemáticos de acordo com a estrutura da fórmula.
        \item Usar recursos da pilha AMS para numeração, alinhamento e operadores.
        \item Criar macros semânticas para notação recorrente.
        \item Configurar cabeçalhos e rodapés de forma controlada.
        \item Refatorar trechos matemáticos frágeis.
    \end{itemize}
\end{frame}

\begin{frame}{Problema do Dia}
    Um documento matemático pode compilar e ainda assim estar mal estruturado.
    \vfill
    \begin{itemize}
        \item Fórmulas longas sem ambiente adequado.
        \item Operadores escritos como texto improvisado.
        \item Delimitadores escolhidos mecanicamente.
        \item Labels e referências difíceis de manter.
        \item Notação repetida sem macros próprias.
    \end{itemize}
\end{frame}

% BLOCO 2 - ESTRUTURA DE PROJETO
\begin{frame}[fragile]{Arquivo Principal}
    O arquivo principal deve mostrar a estrutura do projeto.
\begin{lstlisting}[language={[latex]TeX}]
\documentclass{book}

\usepackage{mathtools}
\usepackage{amssymb}

\begin{document}
\input{introducao}
\input{preliminares}
\input{resultados}
\end{document}
\end{lstlisting}
\end{frame}

\begin{frame}{Convenções de Organização}
    \begin{itemize}
        \item Arquivos com nomes descritivos: \texttt{preliminares.tex}, \texttt{resultados.tex}.
        \item Labels com prefixos: \texttt{sec:}, \texttt{eq:}, \texttt{fig:}, \texttt{tab:}.
        \item Macros no preâmbulo ou em arquivo dedicado quando crescerem.
        \item Cabeçalhos e rodapés devem orientar, não competir com o texto.
        \item Exemplos pequenos devem continuar compiláveis isoladamente.
    \end{itemize}
\end{frame}

\begin{frame}{Política de Página}
    Cabeçalho e rodapé são decisões estruturais do documento.
    \vfill
    \begin{itemize}
        \item O que ajuda o leitor a se localizar?
        \item O documento é frente e verso ou apenas digital?
        \item A abertura de capítulo deve ter o mesmo estilo das páginas comuns?
        \item O número de página deve ficar no rodapé ou no lado externo?
    \end{itemize}
\end{frame}

\begin{frame}[fragile]{Rodapé Simples}
    Para documentos curtos, muitas vezes basta numerar a página.
\begin{lstlisting}[language={[latex]TeX}]
\usepackage{fancyhdr}

\pagestyle{fancy}
\fancyhf{}
\fancyfoot[C]{\thepage}
\renewcommand{\headrulewidth}{0pt}
\renewcommand{\footrulewidth}{0pt}
\end{lstlisting}
\end{frame}

\begin{frame}[fragile]{Cabeçalho em Documento Longo}
\begin{lstlisting}[language={[latex]TeX}]
\usepackage{fancyhdr}

\pagestyle{fancy}
\fancyhf{}
\fancyhead[LE]{\nouppercase{\leftmark}}
\fancyhead[RO]{\nouppercase{\rightmark}}
\fancyfoot[C]{\thepage}
\end{lstlisting}
    \texttt{LE} e \texttt{RO} distinguem páginas pares e ímpares.
\end{frame}

\begin{frame}[fragile]{Marcas de Capítulo e Seção}
\begin{lstlisting}[language={[latex]TeX}]
\fancyhead[LO]{\nouppercase{\rightmark}}
\fancyhead[RE]{\nouppercase{\leftmark}}
\fancyhead[LE,RO]{\thepage}
\end{lstlisting}
    \begin{itemize}
        \item \texttt{\textbackslash leftmark}: normalmente capítulo.
        \item \texttt{\textbackslash rightmark}: normalmente seção.
        \item \texttt{\textbackslash nouppercase}: evita caixa alta automática.
    \end{itemize}
\end{frame}

\begin{frame}[fragile]{Página Especial: \texttt{plain}}
    Páginas de abertura costumam precisar de um estilo mais discreto.
\begin{lstlisting}[language={[latex]TeX}]
\fancypagestyle{plain}{%
  \fancyhf{}
  \fancyfoot[C]{\thepage}
  \renewcommand{\headrulewidth}{0pt}
}
\end{lstlisting}
    \textbf{Critério:} mostre apenas informação útil para navegação.
\end{frame}

\begin{frame}[fragile]{Aviso de \texttt{headheight}}
    Se o cabeçalho não cabe na área reservada, o log pode avisar.
\begin{lstlisting}[language={[latex]TeX}]
\setlength{\headheight}{15pt}
\end{lstlisting}
    \textbf{Interpretação:} o problema não é só estético; o layout da página precisa reservar espaço suficiente.
\end{frame}

% BLOCO 3 - PILHA AMS
\begin{frame}[fragile]{A Pilha AMS}
\begin{lstlisting}[language={[latex]TeX}]
\usepackage{mathtools} % carrega amsmath
\usepackage{amssymb}
\end{lstlisting}
    \begin{itemize}
        \item \textbf{\texttt{amsmath}}: ambientes, numeração, operadores e estruturas matemáticas.
        \item \textbf{\texttt{amssymb}}: símbolos adicionais.
        \item \textbf{\texttt{mathtools}}: extensões práticas sobre \texttt{amsmath}.
    \end{itemize}
\end{frame}

\begin{frame}{Decisão Principal}
    Antes de escolher um ambiente, identifique a estrutura matemática.
    \vfill
    \begin{itemize}
        \item Uma equação isolada?
        \item Uma equação longa com um único número?
        \item Um cálculo com alinhamento vertical?
        \item Um grupo de equações independentes?
        \item Um sistema com numeração subordinada?
    \end{itemize}
\end{frame}

% BLOCO 4 - AMBIENTES
\begin{frame}[fragile]{Equação Isolada}
\begin{lstlisting}[language={[latex]TeX}]
\begin{equation}\label{eq:energia}
E = mc^2
\end{equation}

\[
E = mc^2
\]
\end{lstlisting}
    \textbf{Uso:} uma fórmula com papel próprio no texto.
\end{frame}

\begin{frame}[fragile]{Equação Longa: \texttt{split}}
\begin{lstlisting}[language={[latex]TeX}]
\begin{equation}\label{eq:funcional}
\begin{split}
J(u)
  &= \int_\Omega |\nabla u|^2\,dx \\
  &\quad + \int_\Omega V|u|^2\,dx
       - \int_\Omega F(x,u)\,dx .
\end{split}
\end{equation}
\end{lstlisting}
    \textbf{Ideia:} várias linhas, mas uma única equação numerada.
\end{frame}

\begin{frame}[fragile]{Cálculo Alinhado: \texttt{align}}
\begin{lstlisting}[language={[latex]TeX}]
\begin{align}
\|u+v\|^2
  &= \langle u+v, u+v\rangle \\
  &= \|u\|^2 + 2\langle u,v\rangle + \|v\|^2 .
\end{align}
\end{lstlisting}
    O ponto de alinhamento vem antes do símbolo relacional: \texttt{\&=}.
\end{frame}

\begin{frame}[fragile]{Evite \texttt{eqnarray}}
\begin{lstlisting}[language={[latex]TeX}]
% Evite em documentos novos
\begin{eqnarray}
a &=& b + c
\end{eqnarray}

% Prefira
\begin{align}
a &= b + c
\end{align}
\end{lstlisting}
    \texttt{eqnarray} produz espaçamento inadequado e lida mal com numeração em linhas longas.
\end{frame}

\begin{frame}[fragile]{Grupos Sem Alinhamento: \texttt{gather}}
\begin{lstlisting}[language={[latex]TeX}]
\begin{gather}
A+B+C = 0,\\
\int_\Omega f(x)\,dx = 1.
\end{gather}
\end{lstlisting}
    Use quando as equações pertencem ao mesmo bloco, mas não compartilham uma coluna de alinhamento.
\end{frame}

% BLOCO 5 - NUMERACAO
\begin{frame}[fragile]{Referências de Equações}
\begin{lstlisting}[language={[latex]TeX}]
\begin{equation}\label{eq:identidade}
(a+b)^2 = a^2 + 2ab + b^2
\end{equation}

Pela identidade \eqref{eq:identidade}, temos...
\end{lstlisting}
    \textbf{\texttt{\textbackslash eqref}} imprime os parênteses e preserva o padrão tipográfico.
\end{frame}

\begin{frame}[fragile]{Controlando Números}
\begin{lstlisting}[language={[latex]TeX}]
\begin{align}
f(x) &= x^2 + 2x + 1 \notag\\
     &= (x+1)^2. \label{eq:quadrado}
\end{align}
\end{lstlisting}
    Nem toda linha precisa receber número. Numere apenas o que será referenciado ou tem papel próprio.
\end{frame}

\begin{frame}[fragile]{Numeração Subordinada}
\begin{lstlisting}[language={[latex]TeX}]
\begin{subequations}\label{eq:sistema}
\begin{align}
-\Delta u &= f && \text{em } \Omega,\\
u &= 0 && \text{em } \partial\Omega.
\end{align}
\end{subequations}
\end{lstlisting}
    Útil para sistemas: uma referência para o conjunto e referências para linhas específicas.
\end{frame}

% BLOCO 6 - TEXTO E JUSTIFICATIVAS
\begin{frame}[fragile]{Texto Dentro de Matemática}
\begin{lstlisting}[language={[latex]TeX}]
\[
P_{r-j} =
\begin{cases}
0, & \text{se } r-j \text{ e impar},\\
1, & \text{se } r-j \text{ e par}.
\end{cases}
\]
\end{lstlisting}
    Use \texttt{\textbackslash text} para palavras e frases dentro de fórmulas.
\end{frame}

\begin{frame}[fragile]{Justificativas em Cálculos}
\begin{lstlisting}[language={[latex]TeX}]
\begin{align}
x &= y_1 - y_2 + y_3
     && \text{por \eqref{eq:recorrencia}}\\
  &= y^\prime \circ y^*
     && \text{pela definicao.}
\end{align}
\end{lstlisting}
    A segunda coluna separa a expressão matemática da justificativa textual.
\end{frame}

\begin{frame}[fragile]{Interrupções com Alinhamento}
\begin{lstlisting}[language={[latex]TeX}]
\begin{align}
A_1 &= N_0(\lambda;\Omega') - \phi(\lambda;\Omega'),\\
A_2 &= \phi(\lambda;\Omega') - \phi(\lambda;\Omega),
\shortintertext{e portanto}
A_3 &= N(\lambda;\omega).
\end{align}
\end{lstlisting}
    \texttt{\textbackslash shortintertext} vem de \texttt{mathtools}.
\end{frame}

% BLOCO 7 - OPERADORES E MACROS
\begin{frame}[fragile]{Operadores Declarados}
\begin{lstlisting}[language={[latex]TeX}]
\DeclareMathOperator{\rank}{rank}
\DeclareMathOperator{\tr}{tr}

\[
\rank A + \rank B \geq \rank(A+B),
\qquad
\tr(AB) = \tr(BA).
\]
\end{lstlisting}
    Operadores declarados recebem fonte e espaçamento apropriados.
\end{frame}

\begin{frame}[fragile]{Operadores com Limites}
\begin{lstlisting}[language={[latex]TeX}]
\DeclareMathOperator*{\argmin}{arg\,min}

\[
u = \argmin_{v\in V} J(v).
\]
\end{lstlisting}
    A forma estrelada permite limites abaixo/acima em modo display.
\end{frame}

\begin{frame}[fragile]{Macros Semânticas}
\begin{lstlisting}[language={[latex]TeX}]
\newcommand{\R}{\mathbb{R}}
\newcommand{\eps}{\varepsilon}
\newcommand{\dd}{\,d}

\[
\int_{\R^n} f(x)\dd x < \eps.
\]
\end{lstlisting}
    Uma macro deve expressar notação recorrente, não apenas economizar digitação.
\end{frame}

\begin{frame}[fragile]{Comandos com Argumentos}
\begin{lstlisting}[language={[latex]TeX}]
\newcommand{\norma}[1]{\lVert#1\rVert}
\newcommand{\derivada}[2]{\frac{d #1}{d #2}}

\[
\norma{u}
\qquad
\derivada{y}{x}
\]
\end{lstlisting}
    \texttt{\#1} e \texttt{\#2} são substituídos pelos argumentos do comando.
\end{frame}

\begin{frame}[fragile]{\texttt{\textbackslash newcommand} e \texttt{\textbackslash def}}
\begin{lstlisting}[language={[latex]TeX}]
% Preferivel em documentos LaTeX
\newcommand{\R}{\mathbb{R}}

% Primitivo TeX: comum em codigo interno
\def\R{\mathbb{R}}
\end{lstlisting}
    \texttt{\textbackslash newcommand} verifica conflitos. \texttt{\textbackslash def} sobrescreve de forma mais direta.
\end{frame}

\begin{frame}[fragile]{Redefinição Deliberada}
\begin{lstlisting}[language={[latex]TeX}]
\renewcommand{\epsilon}{\varepsilon}
\providecommand{\dd}{\,d}
\end{lstlisting}
    \begin{itemize}
        \item \texttt{\textbackslash renewcommand}: redefinição explícita.
        \item \texttt{\textbackslash providecommand}: define apenas se ainda não existir.
    \end{itemize}
\end{frame}

% BLOCO 8 - DELIMITADORES
\begin{frame}[fragile]{Delimitadores Pareados}
\begin{lstlisting}[language={[latex]TeX}]
\DeclarePairedDelimiter{\abs}{\lvert}{\rvert}
\DeclarePairedDelimiter{\norm}{\lVert}{\rVert}

\[
\abs{x}
\qquad
\abs*{\frac{x}{1+x^2}}
\qquad
\norm[\Big]{\sum_{k=1}^n v_k}
\]
\end{lstlisting}
    \texttt{mathtools} oferece forma natural, automática e tamanho explícito.
\end{frame}

\begin{frame}[fragile]{Produto Interno}
\begin{lstlisting}[language={[latex]TeX}]
\DeclarePairedDelimiterX{\inner}[2]
  {\langle}{\rangle}{#1,#2}

\[
\inner{u}{v}
\qquad
\inner*{u}{\frac{v}{\|v\|}}
\]
\end{lstlisting}
    A macro concentra a decisão tipográfica da notação.
\end{frame}

% BLOCO 9 - MATRIZES E CASOS
\begin{frame}[fragile]{Matrizes}
\begin{lstlisting}[language={[latex]TeX}]
\[
A =
\begin{pmatrix}
1 & 0\\
0 & 1
\end{pmatrix},
\qquad
\det A =
\begin{vmatrix}
a & b\\
c & d
\end{vmatrix}.
\]
\end{lstlisting}
\end{frame}

\begin{frame}[fragile]{Casos com \texttt{dcases}}
\begin{lstlisting}[language={[latex]TeX}]
\[
f(x) =
\begin{dcases}
\frac{\sin x}{x}, & x\neq 0,\\
1, & x = 0.
\end{dcases}
\]
\end{lstlisting}
    \texttt{dcases}, de \texttt{mathtools}, usa estilo de display nos casos.
\end{frame}

\begin{frame}[fragile]{Reticências Semânticas}
\begin{lstlisting}[language={[latex]TeX}]
\[
a_1, a_2, \dotsc, a_n
\qquad
a_1 + a_2 + \dotsb + a_n
\qquad
a_1 a_2 \dotsm a_n
\]
\end{lstlisting}
    O comando informa o contexto: vírgulas, operadores binários ou multiplicação.
\end{frame}

% BLOCO 10 - LABORATORIO
\begin{frame}[fragile]{Laboratório: Refatoração Matemática}
    \textbf{Objetivo:} transformar um trecho que apenas compila em um trecho matematicamente estruturado.

    \textbf{Arquivo:} \texttt{ams-refatoracao.tex}

\begin{lstlisting}[language={[latex]TeX}]
\begin{eqnarray}
||u+v||^2 &=& <u+v,u+v>\\
&=& ||u||^2 + 2<u,v> + ||v||^2.
\end{eqnarray}
\end{lstlisting}
\end{frame}

\begin{frame}{Laboratório: Tarefas}
    \begin{enumerate}
        \item Carregar \texttt{mathtools} e \texttt{amssymb}.
        \item Substituir \texttt{eqnarray} por \texttt{align}.
        \item Criar macros para norma e produto interno.
        \item Numerar apenas a linha final.
        \item Referenciar a identidade com \texttt{\textbackslash eqref}.
    \end{enumerate}
    \vfill
    \textbf{Observe:} a mudança melhora o código-fonte e o PDF.
\end{frame}

\begin{frame}[fragile]{Laboratório: Uma Solução}
\begin{lstlisting}[language={[latex]TeX}]
\DeclarePairedDelimiter{\norm}{\lVert}{\rVert}
\DeclarePairedDelimiterX{\inner}[2]{\langle}{\rangle}{#1,#2}

\begin{align}
\norm{u+v}^2
  &= \inner{u+v}{u+v} \notag\\
  &= \norm{u}^2 + 2\inner{u}{v} + \norm{v}^2.
     \label{eq:paralelogramo}
\end{align}
\end{lstlisting}
\end{frame}

\begin{frame}[fragile]{Laboratório: Cabeçalho e Rodapé}
    \textbf{Objetivo:} configurar um estilo de página para um documento técnico.

    \textbf{Tarefa:}
    \begin{enumerate}
        \item Criar um documento \texttt{book} com duas seções.
        \item Carregar \texttt{fancyhdr}.
        \item Colocar número de página no lado externo.
        \item Usar capítulo/seção no cabeçalho.
    \end{enumerate}
\begin{lstlisting}[language={[latex]TeX}]
\documentclass[twoside]{book}
\usepackage{fancyhdr}
\end{lstlisting}
\end{frame}

\begin{frame}[fragile]{Laboratório: Estilo de Página}
\begin{lstlisting}[language={[latex]TeX}]
\pagestyle{fancy}
\fancyhf{}
\fancyhead[LE,RO]{\thepage}
\fancyhead[LO]{\nouppercase{\rightmark}}
\fancyhead[RE]{\nouppercase{\leftmark}}

\fancypagestyle{plain}{%
  \fancyhf{}
  \fancyfoot[C]{\thepage}
}
\end{lstlisting}
    \textbf{Observe:} páginas comuns e páginas de abertura não precisam ter o mesmo desenho.
\end{frame}

% BLOCO 11 - FECHAMENTO
\begin{frame}{Recapitulação}
    \begin{itemize}
        \item A escolha do ambiente deve seguir a estrutura da fórmula.
        \item \texttt{amsmath} resolve alinhamento, numeração e operadores.
        \item \texttt{mathtools} amplia delimitadores e construções práticas.
        \item Macros semânticas tornam a notação consistente.
        \item Cabeçalhos e rodapés são parte da arquitetura do documento.
        \item Refatorar matemática é melhorar manutenção, não apenas aparência.
    \end{itemize}
\end{frame}

\begin{frame}{Próximo Dia}
    \begin{itemize}
        \item Estrutura de apresentações com Beamer.
        \item Frames, blocos e colunas.
        \item Overlays e progressão de aula.
        \item Código em slides.
        \item Boas práticas para apresentação técnica.
    \end{itemize}
\end{frame}

\end{document}
